\chapter{Introduction}
\section{Overview}
In this template, we refer to several sources such as \cite{sample_article}, \cite{sample_book}, and the BibLaTeX documentation \cite{sample_online}.
\section{Motivation}
\section{Objective}
\section{Scope}
\section{SDG Justification}
\section{Existing System}
\section{Proposed System}
\section{Work Plan}

B.M.S. College of Engineering (BMSCE), founded in 1946 by Late Sri B. M. Sreenivasaiah, 
is the first private engineering college in India. Located in the heart of Bengaluru, 
Karnataka, the institution has grown to become one of the most prestigious engineering 
colleges in the country. 

BMSCE is autonomous, affiliated with Visvesvaraya Technological University (VTU), Belagavi, 
and is recognized by the All India Council for Technical Education (AICTE). The college is 
accredited by the National Board of Accreditation (NBA) and the National Assessment and 
Accreditation Council (NAAC) with an “A++” grade.

The institution offers 14 undergraduate and 15 postgraduate programs across various 
engineering disciplines, management, and computer applications. It also has strong research 
centers in multiple domains, encouraging innovation and industry collaboration. 

BMSCE is known for its state-of-the-art infrastructure, highly qualified faculty, and vibrant 
student community. The college emphasizes holistic development through academics, research, 
co-curricular and extracurricular activities, and has a strong alumni network spread 
globally. 

With its motto *“Service to Humanity is Service to God”*, BMSCE continues to foster 
excellence in education, research, and innovation.
